% ----------------------
\subsection{Section 68 (1 November 1831)}
% ----------------------
	\label{DC68}
	
	% -----
	\subsubsection{Relation to section 1}
	% -----
	
		Section 68 was given at the same conference during which JS received \hyperref[DC1]{section 1}.
	
	% -----
	\subsubsection{Historical Background}
	% -----
		\label{DC68-HB}
		
		% --
		\paragraph{JSP}
		% --
		
			``On 1--2 November 1831, ten elders convened a conference in Hiram, Ohio, to discuss the publication of the Book of Commandments, a compilation of JS's revelations. According to a later JS history, four of the conference attendees---Orson Hyde, Luke Johnson, Lyman Johnson, and William E. McLellin---approached JS during the conference and requested to know the Lord's will concerning them. This revelation came in response to their inquiry. The revelation provided more information about the evangelizing duties of the four men specifically and of elders in general. While Hyde, McLellin, and Luke Johnson were all ordained to the high priesthood at a conference held in Orange, Cuyahoga County, Ohio, a week earlier, Lyman Johnson was ordained to the high priesthood at the Hiram conference on 2 November.
			
			``After closing the portion of the revelation addressed specifically to the four men with an `Amen,' the document shifts its audience to the church in general and gives additional information about the office of bishop, as well as counsel to members of the church `in Zion' about teaching and baptizing their children and avoiding idleness and greed. The text may originally have been dictated as two discrete revelations, which, like some other revelations closely related in time or content, were then copied together and presented as a single, unified text. All extant copies of the text---whether in manuscript or published form---present both parts as one revelation.
			
			``The original manuscript of the revelation is not extant, and the conference minutes do not mention the revelation. However, the copy in Revelation Book 1 is dated 1 November 1831 and a heading states that it was `given in Hiram Nov. 1. 1831.' John Whitmer and Oliver Cowdery copied the revelation into Revelation Book 1, probably soon after its dictation.''
			
		% --
		\paragraph{HDDC} 
		% --
		
			``There are several rather important alterations in the text of this revelation. Most of them are connected with explaining the duties and qualifications of a Bishop, and the responsibility of the person who calls another to that office in the priesthood. In 1831, when this revelation was given, there was no First Presidency of the Church, and it was not until March 1835 that much of the information about bishops was revealed (see \hyperref[DC107]{DC 107}). Therefore, in light of these later revelations, Section 68 contained several changes when it was published in the Kirtland reprint of the \uline{Evening and Morning Star}, June 1835, and these changes have been retained in all printings since then. Most of the changes are identified in the Text Analysis by the letters E, F, G, and I on page 111 and letters A and B on page 112.
			
			``Other changes of some importance are `D' on page 111 and `D--H' on page 112 of the Text Analysis'' (856--857).%
				\footnote{%
					% \label{}%
					See HDDC 860--865 for Woodford's analysis. The bulk of the changes are the addition of verses 15--21 to the section, almost all of which was not in the original version. See below.
					}
			
	% -----
	\subsubsection{Versions}
	% -----
		\label{DC68-V}
		
		See the \href{https://drive.google.com/file/d/1goAvNz8jS4R8slYfMG_db55Ty2z_vmgK/view?usp=sharing}{sections comparison} document.
	
		\begin{compactdesc}
			\item[JSP] \hyperref[EMS-1831-JS-1-NOV-1831-A]{1 November 1831}
			\item[EMS] 1:5 (October 1832), 35 (verses 1--15a and 22--35 only)
			\item[BC] ---
			\item[EMS-R] 1:15 (June 1835), 73--74\footnotemark
			\item[DC 1835] Section 22, 147--149
			\item[TS] 5:8 (15 April 1844), 496--497
			\item[DC 1844] Section 22, 214--217
			\item[MS] 5:12 (May 1845), 186--187
		\end{compactdesc}
			\footnotetext{%
				% \label{}%
				See the historical background info from HDDC above for more here. If you go check my copy of the reprint, it's actually 1:5, reprinting October 1832, page 73.
				}

	% -----
	\subsubsection{Section 68:1} % *DC68-1 
	% -----
		\label{DC68-1}

		MY servant, Orson Hyde, was called by his ordination to proclaim the everlasting gospel, by the Spirit of the living God, from people to people, and from land to land, in the congregations of the wicked, in their synagogues, reasoning with and expounding all scriptures unto them.

		% \begin{framed}
		% \scriptsize
			% \begin{compactdesc}
				% \item[JSP] 
				% % \item[BC] 
				% % \item[]
			% \end{compactdesc}
		% \end{framed}

	% -----
	\subsubsection{Section 68:2} % *DC68-2 
	% -----
		\label{DC68-2}

		And, behold, and lo, this is an ensample unto all those who were ordained unto this priesthood,%
			\footnote{%
				% \label{}%
				Could be referring to the ``high priesthood,'' which was introduced shortly before this.
				}
		whose mission is appointed unto them to go forth---

		% \begin{framed}
		% \scriptsize
			% \begin{compactdesc}
				% \item[JSP] 
				% % \item[BC] 
				% % \item[]
			% \end{compactdesc}
		% \end{framed}

	% -----
	\subsubsection{Section 68:3} % *DC68-3 
	% -----
		\label{DC68-3}

		And this is the ensample unto them, that they shall speak as they are moved upon by the Holy Ghost.

		% \begin{framed}
		% \scriptsize
			% \begin{compactdesc}
				% \item[JSP] 
				% % \item[BC] 
				% % \item[]
			% \end{compactdesc}
		% \end{framed}

	% -----
	\subsubsection{Section 68:4} % *DC68-4 
	% -----
		\label{DC68-4}

		And whatsoever they shall speak when moved upon by the Holy Ghost shall be scripture, shall be the will of the Lord, shall be the mind of the Lord, shall be the word of the Lord, shall be the voice of the Lord, and the power of God unto salvation.%
			\footnote{%
				% \label{}%
				I almost want to restrict this scripture to the missionaries that the Lord is specifically referring to here; otherwise I'm not totally sure what to do with this idea.
				}

		% \begin{framed}
		% \scriptsize
			% \begin{compactdesc}
				% \item[JSP] 
				% % \item[BC] 
				% % \item[]
			% \end{compactdesc}
		% \end{framed}

	% -----
	\subsubsection{Section 68:5} % *DC68-5 
	% -----
		\label{DC68-5}

		Behold, this is the promise of the Lord unto you, O ye my servants.

		% \begin{framed}
		% \scriptsize
			% \begin{compactdesc}
				% \item[JSP] 
				% % \item[BC] 
				% % \item[]
			% \end{compactdesc}
		% \end{framed}

	% -----
	\subsubsection{Section 68:6} % *DC68-6 
	% -----
		\label{DC68-6}

		Wherefore, be of good cheer, and do not fear, for I the Lord am with you, and will stand by you; and ye shall bear record of me, even Jesus Christ, that I am the Son of the living God, that I was, that I am, and that I am to come.%
			\footnote{%
				% \label{}%
				For some reason, this reminds me of \hyperref[EXODUS-3-14]{Exodus 3:14}, but also stuff like \hyperref[DC93-24]{DC 93:24}.
				}

		% \begin{framed}
		% \scriptsize
			% \begin{compactdesc}
				% \item[JSP] 
				% % \item[BC] 
				% % \item[]
			% \end{compactdesc}
		% \end{framed}

	% -----
	\subsubsection{Section 68:7} % *DC68-7 
	% -----
		\label{DC68-7}

		This is the word of the Lord unto you, my servant Orson Hyde, and also unto my servant Luke Johnson, and unto my servant Lyman Johnson, and unto my servant William E. McLellin, and unto all the faithful elders of my church---

		% \begin{framed}
		% \scriptsize
			% \begin{compactdesc}
				% \item[JSP] 
				% % \item[BC] 
				% % \item[]
			% \end{compactdesc}
		% \end{framed}

	% -----
	\subsubsection{Section 68:8} % *DC68-8 
	% -----
		\label{DC68-8}

			\footnote{%
				% \label{}%
				This verse and the next are drawing from a combination of \hyperref[MARK-16-15]{Mark 16:15--16} and \hyperref[MATTHEW-28-19]{Matthew 28:19}.
				}%
		Go ye into all the world, preach the gospel to every creature, acting in the authority%
		\footnote{%
			% \label{}%
			Once again, I think it's referring to the ``high priesthood.''
			} 
		which I have given you, baptizing in the name of the Father, and of the Son, and of the Holy Ghost.

		% \begin{framed}
		% \scriptsize
			% \begin{compactdesc}
				% \item[JSP] 
				% % \item[BC] 
				% % \item[]
			% \end{compactdesc}
		% \end{framed}

	% -----
	\subsubsection{Section 68:9} % *DC68-9 
	% -----
		\label{DC68-9}

		And he that believeth and is baptized shall be saved, and he that believeth not shall be damned.

		% \begin{framed}
		% \scriptsize
			% \begin{compactdesc}
				% \item[JSP] 
				% % \item[BC] 
				% % \item[]
			% \end{compactdesc}
		% \end{framed}

	% -----
	\subsubsection{Section 68:10} % *DC68-10 
	% -----
		\label{DC68-10}

		And he that believeth shall be blest with signs following, even as it is written.

		% \begin{framed}
		% \scriptsize
			% \begin{compactdesc}
				% \item[JSP] 
				% % \item[BC] 
				% % \item[]
			% \end{compactdesc}
		% \end{framed}

	% -----
	\subsubsection{Section 68:11} % *DC68-11 
	% -----
		\label{DC68-11}

		And unto you it shall be given to know the signs of the times, and the signs of the coming of the Son of Man;

		% \begin{framed}
		% \scriptsize
			% \begin{compactdesc}
				% \item[JSP] 
				% % \item[BC] 
				% % \item[]
			% \end{compactdesc}
		% \end{framed}

	% -----
	\subsubsection{Section 68:12} % *DC68-12 
	% -----
		\label{DC68-12}

		And of as many as the Father shall bear record, to you shall be given power%
			\footnote{%
				% \label{}%
				Note that this ``power'' is missing from the JSP version.
				}
		to seal them up unto eternal life.  Amen.

		\begin{framed}
		\scriptsize
			\begin{compactdesc}
				\item[JSP] \& of as many as the Father shall bear record to you it shall be given to seal them up unto Eternal life Amen---
				\item[EMS] and of as many as the Father shall be record, to you it shall be given power to seal them up unto eternal life: Amen.
				\item[EMS-R] and of as many as the Father shall bear record, to you it shall be given power to seal them up unto eternal life: Amen.
				\item[DC 1835] and of as many as the Father shall bear record, to you it shall be given power to seal them up unto eternal life: Amen.
				% \item[BC] 
				% \item[]
			\end{compactdesc}
		\end{framed}

	% -----
	\subsubsection{Section 68:13} % *DC68-13 
	% -----
		\label{DC68-13}

		And now, concerning the items in addition to the covenants and commandments, they are these---

		\begin{framed}
		\scriptsize
			\begin{compactdesc}
				\item[JSP] And now conc[e]rning the items in addition to the Laws \& commandments they are these
				\item[EMS] And now, concerning the items in addition to the Laws and commandments, they are these:
				\item[EMS-R] And now concerning the items in addition to the covenants and commandments, they are these:
				\item[DC 1835] 2 And now concerning the items in addition to the covenants and commandments, they are these:
				% \item[BC] 
				% \item[]
			\end{compactdesc}
		\end{framed}

	% -----
	\subsubsection{Section 68:14} % *DC68-14 
	% -----
		\label{DC68-14}

			\footnote{%
				% \label{}%
				Here begins a seemingly important but very confusing set of passages on the Aaronic priesthood. I'm not sure how much sense this really makes. Note that after this verse, very little of what follows is in the original version; see the \hyperref[DC68-HB]{historical background} above, especially HDDC---this material gets added later, after \hyperref[DC107]{DC 107} comes along. So this stuff is an explication or expansion of the things in DC 107.
				}%
		There remain hereafter, in the due time of the Lord, other bishops to be set apart unto the church, to minister even according to the first;

		\begin{framed}
		\scriptsize
			\begin{compactdesc}
				\item[JSP] there rema[i]neth hereafter in the due time of the Lord other Bishops to be set apart unto the church to minister even according to the first
				\item[EMS] There remaineth hereafter in the due time of the Lord, other bishops to be set apart unto the church, to minister even according to the first; 
				\item[EMS-R] there remaineth hereafter, in the due time of the Lord, other bishops to be set apart unto the church to minister even according to the first:
				\item[DC 1835] There remaineth hereafter in the due time of the Lord, other bishops to be set apart unto the church to minister even according to the first:
				% \item[BC] 
				% \item[]
			\end{compactdesc}
		\end{framed}

	% -----
	\subsubsection{Section 68:15} % *DC68-15 
	% -----
		\label{DC68-15}

		Wherefore they shall be high priests who are worthy, and they shall be appointed by the First Presidency of the Melchizedek Priesthood,%
			\footnote{%
				% \label{}%
				Note the differences in the older versions; see also the \hyperref[DC72-HB]{historical background} to section 72, which talks about the conference they ended up having.
				} 
		except they be literal descendants of Aaron.

		\begin{framed}
		\scriptsize
			\begin{compactdesc}
				\item[JSP] wherefore it shall be an high priest who is worthy \& he shall be appointed by a confrenc of high priests
				\item[EMS] wherefore it shall be an high priest who is worthy; and he shall be appointed by a conference of high priests.
				\item[EMS-R] Wherefore they shall be high priests who are worthy, and they shall be appointed by the first presidency of the Melchisedek priesthood, except they be literal descendants of Aaron;
				\item[DC 1835] wherefore they shall be high priests who are worthy, and they shall be appointed by the first presidency of the Melchizedek priesthood, except they be literal descendants of Aaron,
				% \item[BC] 
				% \item[]
			\end{compactdesc}
		\end{framed}

	% -----
	\subsubsection{Section 68:16} % *DC68-16 
	% -----
		\label{DC68-16}

		And if they be literal descendants of Aaron they have a legal right to the bishopric,%
			\footnote{%
				% \label{}%
				Reminds me of \hyperref[HEBREWS-5-4]{Hebrews 5:4}.
				}
		if they are the firstborn among the sons of Aaron;%
			\footnote{%
				% \label{}%
				A lot of this stuff in these passages just doesn't make a lot of sense to me. How do you know if you're a ``literal'' descendent of Aaron? And what does it mean to be the ``firstborn'' in this sense? (You're going to have to look to \hyperref[DC68-21]{verse 21} here.) The relationship to the ``high priesthood'' doesn't make a lot of sense to me either.
				}

		\begin{framed}
		\scriptsize
			\begin{compactdesc}
				\item[JSP] ---
				\item[EMS] ---
				\item[EMS-R] and if they be literal descendants of Aaron, they have a legal right to the bishopric, if they are the first born among the sons of Aaron:
				\item[DC 1835] and if they be literal descendants of Aaron, they have a legal right to the bishopric, if they are the first born among the sons of Aaron:
				% \item[BC] 
				% \item[]
			\end{compactdesc}
		\end{framed}

	% -----
	\subsubsection{Section 68:17} % *DC68-17 
	% -----
		\label{DC68-17}

		For the firstborn holds the right of the presidency over this priesthood, and the keys or authority of the same.%
			\footnote{%
				% \label{}%
				Note the familial descendance of the priesthood here, and we're only speaking of the Aaronic priesthood here.
				}

		\begin{framed}
		\scriptsize
			\begin{compactdesc}
				\item[JSP] ---
				\item[EMS] ---
				\item[EMS-R] for the first born holds the right of presidency over this priesthood, and the keys or authority of the same.
				\item[DC 1835] for the first born holds the right of presidency over this priesthood, and the keys or authority of the same.
				% \item[BC] 
				% \item[]
			\end{compactdesc}
		\end{framed}

	% -----
	\subsubsection{Section 68:18} % *DC68-18 
	% -----
		\label{DC68-18}

		No man has a legal right to this office, to hold the keys of this priesthood, except he be a literal descendant and the firstborn of Aaron.

		\begin{framed}
		\scriptsize
			\begin{compactdesc}
				\item[JSP] ---
				\item[EMS] ---
				\item[EMS-R] No man has a legal right to this office, to hold the keys of this priesthood, except he be a literal descendant and the first born of Aaron:
				\item[DC 1835] No man has a legal right to this office, to hold the keys of this priesthood, except he be a literal descendant and the first born of Aaron:
				% \item[BC] 
				% \item[]
			\end{compactdesc}
		\end{framed}

	% -----
	\subsubsection{Section 68:19} % *DC68-19 
	% -----
		\label{DC68-19}

		But, as a high priest of the Melchizedek Priesthood has authority to officiate in all the lesser offices he may officiate in the office of bishop when no literal descendant of Aaron can be found, provided he is called and set apart and ordained unto this power, under the hands of the First Presidency of the Melchizedek Priesthood.

		\begin{framed}
		\scriptsize
			\begin{compactdesc}
				\item[JSP] ---
				\item[EMS] ---
				\item[EMS-R] but as a high priest of the Melchizedek priesthood, has authority to officiate in all the lesser offices, he may officiate in the office of bishop when no literal descendant of Aaron can be found; provided he is called and set apart, and ordained unto this power under the hands of the first presidency of the Melchizedek priesthood.
				\item[DC 1835] but as a high priest of the Melchizedek priesthood, has authority to officiate in all the lesser offices, he may officiate in the office of bishop when no literal descendant of Aaron can be found; provided he is called and set apart, and ordained unto this power under the hands of the first presidency of the Melchizedek priesthood.
				% \item[BC] 
				% \item[]
			\end{compactdesc}
		\end{framed}

	% -----
	\subsubsection{Section 68:20} % *DC68-20 
	% -----
		\label{DC68-20}

		And a literal descendant of Aaron, also, must be designated by this Presidency, and found worthy, and anointed, and ordained under the hands of this Presidency, otherwise they are not legally authorized to officiate in their priesthood.

		\begin{framed}
		\scriptsize
			\begin{compactdesc}
				\item[JSP] ---
				\item[EMS] ---
				\item[EMS-R] And a literal descendant of Aaron, also, must be designated by this presidency, and found worthy, and anointed, and ordained under the hands of this presidency, otherwise they are not legally authorized to officiate in their priesthood:
				\item[DC 1835] And a literal descendant of Aaron, also, must be designated by this presidency, and found worthy, and anointed, and ordained under the hands of this presidency, otherwise they are not legally authorized to officiate in their priesthood:
				% \item[BC] 
				% \item[]
			\end{compactdesc}
		\end{framed}

	% -----
	\subsubsection{Section 68:21} % *DC68-21 
	% -----
		\label{DC68-21}

		But, by virtue of the decree concerning their right of the priesthood descending from father to son, they may claim their anointing if at any time they can prove their lineage, or do ascertain it by revelation%
			\footnote{%
				% \label{}%
				Seems to be referring to patriarchal blessings, but it also talks about being a ``literal descendant'' (verse 18/20), and if that's what we're talking about here, then I don't see how different people in different families could have different lineages. If that's possible then we aren't talking about ``literal descendants.''
				}
		from the Lord under the hands of the above named Presidency.

		\begin{framed}
		\scriptsize
			\begin{compactdesc}
				\item[JSP] ---
				\item[EMS] ---
				\item[EMS-R] but by virtue of the decree concerning their right of the priesthood descending from father to son, they may claim their annointing, if at any time they can prove their lineage, or do ascertain it by revelation from the Lord under the hands of the above named presidency.
				\item[DC 1835] but by virtue of the decree concerning their right of the priesthood descending from father to son, they may claim their annointing, if at any time they can prove their lineage, or do ascertain it by revelation from the Lord under the hands of the above named presidency.
				% \item[BC] 
				% \item[]
			\end{compactdesc}
		\end{framed}

	% -----
	\subsubsection{Section 68:22} % *DC68-22 
	% -----
		\label{DC68-22}

		And again, no bishop or high priest who shall be set apart for this ministry shall be tried or condemned for any crime, save it be before the First Presidency of the church;

		\begin{framed}
		\scriptsize
			\begin{compactdesc}
				\item[JSP] And again no Bishop or judge which shall be set apart for this ministry shall be tried or condemned for any crime save it be before a confrence of high priests
				\item[EMS] And again, no bishop or judge, which shall be set apart for this ministry, shall be tried or condemned for any crime, save it be before a conference of high priests;
				\item[EMS-R] And again, no bishop or high priest, who shall be set apart for this ministry, shall be tried or condemned for any crime save it be before the first presidency of the church;
				\item[DC 1835] 3 And again, no bishop or high priest, who shall be set apart for this ministry, shall be tried or condemned for any crime save it be before the first presidency of the church;
				% \item[BC] 
				% \item[]
			\end{compactdesc}
		\end{framed}

	% -----
	\subsubsection{Section 68:23} % *DC68-23 
	% -----
		\label{DC68-23}

		And inasmuch as he is found guilty before this Presidency, by testimony that cannot be impeached, he shall be condemned;

		\begin{framed}
		\scriptsize
			\begin{compactdesc}
				\item[JSP] \& inasmuch as he is found guilty before a confrenc of high priests by testimony that cannot be impeached he shall be condemned
				\item[EMS] and in as much as he is found guilty before a conference of high priests, by testimony that cannot be impeached, he shall be condemned
				\item[EMS-R] and inasmuch as he is found guilty before this presidency, by testimony that cannot be impeached, he shall be condemned, 
				\item[DC 1835] and inasmuch as he is found guilty before this presidency, by testimony that cannot be impeached, he shall be condemned,
				% \item[BC] 
				% \item[]
			\end{compactdesc}
		\end{framed}

	% -----
	\subsubsection{Section 68:24} % *DC68-24 
	% -----
		\label{DC68-24}

		And if he repent he shall be forgiven, according to the covenants and commandments of the church.

		\begin{framed}
		\scriptsize
			\begin{compactdesc}
				\item[JSP] or forgiven according to the Laws of the church
				\item[EMS] or forgiven, according to the laws of the church.
				\item[EMS-R] and if he repents he shall be forgiven, according to the covenants and commandments of the church.
				\item[DC 1835] and if he repents he shall be forgiven, according to the covenants and commandments of the church.
				% \item[BC] 
				% \item[]
			\end{compactdesc}
		\end{framed}

	% -----
	\subsubsection{Section 68:25} % *DC68-25 
	% -----
		\label{DC68-25}

		And again, inasmuch as parents have children in Zion, or in any of her stakes which are organized,%
			\footnote{%
				% \label{}%
				This ``stakes'' bit gets added later; see the older versions.
				}
		that teach%
			\footnote{%
				% \label{}%
				See \hyperref[DEUTERONOMY-6-7]{Deuteronomy 6:7}.
				}
		them not to understand%
			\footnote{%
				% \label{}%
				You have to not only teach it, but teach it in a such a way as to have them \textit{understand}; you could also read it as saying, you have to \textit{teach them the skill of understanding} itself. I like that second reading. See especially \hyperref[MATTHEW-13-19]{Matthew 13:19, 23} and \hyperref[MATTHEW-13-51]{51}, and also \hyperref[2NEPHI-4-5]{2 Nephi 4:5}, \hyperref[DC84-117]{DC 84:117},
				}
		the doctrine of repentance, faith in Christ the Son of the living God, and of baptism and the gift of the Holy Ghost by the laying on of the hands,%
			\footnote{%
				% \label{}%
				An early reference to the ideas in the \hyperref[ARTICLESOFFAITH-4]{fourth article of faith}.
				}
		when eight years old, the sin be upon the heads of the parents.

		\begin{framed}
		\scriptsize
			\begin{compactdesc}
				\item[JSP] And again inasmuch as parents have children in Zion that teach them not to understand the doctrine of repentance faith in Christ the Son of the living God \& of baptism \& the gift of the Holy Spirit by the laying <on> of the hands when eight years old the sin be upon the head of the parents
				\item[EMS] And again, in as much as parents have children in Zion, that teach them not to understand the doctrine of repentance; faith in Christ the Son of the living God; and of baptism and the gift of the Holy Ghost by the laying on of the hands, when eight years old: the sin be upon the head of the parents,
				\item[EMS-R] And again, inasmuch as parents have children in Zion, or in any of her stakes which are organized, that teach them not to understand the doctrine of repentance; faith in Christ the Son of the living God; and of baptism and the gift of the Holy Ghost by the laying on of the hands, when eight years old: the sin be upon the head of the parents,
				\item[DC 1835] 4 And again, inasmuch as parents have children in Zion, or in any of her stakes which are organized, that teach them not to understand the doctrine of repentance; faith in Christ the Son of the living God; and of baptism and the gift of the Holy Ghost by the laying on of the hands, when eight years old, the sin be upon the head of the parents,
				% \item[BC] 
				% \item[]
			\end{compactdesc}
		\end{framed}

	% -----
	\subsubsection{Section 68:26} % *DC68-26 
	% -----
		\label{DC68-26}

		For this shall be a law unto the inhabitants of Zion, or in any of her stakes which are organized.

		\begin{framed}
		\scriptsize
			\begin{compactdesc}
				\item[JSP] for this shall be a Law unto the inhabitants of Zion
				% \item[BC] 
				% \item[]
			\end{compactdesc}
		\end{framed}

	% -----
	\subsubsection{Section 68:27} % *DC68-27 
	% -----
		\label{DC68-27}

		And their children shall be baptized for the remission of their sins when eight years old, and receive the laying on of the hands.

		% \begin{framed}
		% \scriptsize
			% \begin{compactdesc}
				% \item[JSP] 
				% % \item[BC] 
				% % \item[]
			% \end{compactdesc}
		% \end{framed}

	% -----
	\subsubsection{Section 68:28} % *DC68-28 
	% -----
		\label{DC68-28}

		And they shall also teach their children to pray, and to walk uprightly before the Lord.

		% \begin{framed}
		% \scriptsize
			% \begin{compactdesc}
				% \item[JSP] 
				% % \item[BC] 
				% % \item[]
			% \end{compactdesc}
		% \end{framed}

	% -----
	\subsubsection{Section 68:29} % *DC68-29 
	% -----
		\label{DC68-29}

		And the inhabitants of Zion shall also observe the Sabbath day to keep it holy.

		% \begin{framed}
		% \scriptsize
			% \begin{compactdesc}
				% \item[JSP] 
				% % \item[BC] 
				% % \item[]
			% \end{compactdesc}
		% \end{framed}

	% -----
	\subsubsection{Section 68:30} % *DC68-30 
	% -----
		\label{DC68-30}

		And the inhabitants of Zion also shall remember their labors, inasmuch as they are appointed to labor, in all faithfulness; for the idler shall be had in remembrance before the Lord.

		% \begin{framed}
		% \scriptsize
			% \begin{compactdesc}
				% \item[JSP] 
				% % \item[BC] 
				% % \item[]
			% \end{compactdesc}
		% \end{framed}

	% -----
	\subsubsection{Section 68:31} % *DC68-31 
	% -----
		\label{DC68-31}

		Now, I, the Lord, am not well pleased with the inhabitants of Zion, for there are idlers among them; and their children are also growing up in wickedness; they also seek not earnestly the riches of eternity, but their eyes are full of greediness.

		% \begin{framed}
		% \scriptsize
			% \begin{compactdesc}
				% \item[JSP] 
				% % \item[BC] 
				% % \item[]
			% \end{compactdesc}
		% \end{framed}

	% -----
	\subsubsection{Section 68:32} % *DC68-32 
	% -----
		\label{DC68-32}

		These things ought not to be, and must be done away from among them; wherefore, let my servant Oliver Cowdery carry these sayings unto the land of Zion.

		% \begin{framed}
		% \scriptsize
			% \begin{compactdesc}
				% \item[JSP] 
				% % \item[BC] 
				% % \item[]
			% \end{compactdesc}
		% \end{framed}

	% -----
	\subsubsection{Section 68:33} % *DC68-33 
	% -----
		\label{DC68-33}

		And a commandment I give unto them---that he that observeth not his prayers before the Lord in the season thereof, let him be had in remembrance before the judge of my people.

		% \begin{framed}
		% \scriptsize
			% \begin{compactdesc}
				% \item[JSP] 
				% % \item[BC] 
				% % \item[]
			% \end{compactdesc}
		% \end{framed}

	% -----
	\subsubsection{Section 68:34} % *DC68-34 
	% -----
		\label{DC68-34}

		These sayings are true and faithful;%
		\footnote{%
			% \label{}%
			See \hyperref[DC1-37]{DC 1:37}, given at the same time as this section.
			} 
		wherefore, transgress them not, neither take therefrom.

		% \begin{framed}
		% \scriptsize
			% \begin{compactdesc}
				% \item[JSP] 
				% % \item[BC] 
				% % \item[]
			% \end{compactdesc}
		% \end{framed}

	% -----
	\subsubsection{Section 68:35} % *DC68-35 
	% -----
		\label{DC68-35}

		Behold, I am Alpha and Omega, and I come quickly.  Amen.

		% \begin{framed}
		% \scriptsize
			% \begin{compactdesc}
				% \item[JSP] 
				% % \item[BC] 
				% % \item[]
			% \end{compactdesc}
		% \end{framed}